\section{Weak curvature sources and the limits of first-order control}
\label{sec:analytic-limitations}

This section separates three kinds of information: a small Ricci residual,
weak control of the divergence of curvature, and strong control of a
geometrically selected test field.  The examples below show that none of
the relevant transitions follows from a formal integration by parts or
from the displayed first-order norms alone.  The positive results identify
precisely what controlled coordinate geometry supplies.

Throughout this section, two-forms have their exterior-algebra norm and
$W^\pm$ have their Hilbert--Schmidt endomorphism norm.  The norm of $\Rm$
is the full four-index tensor norm.  Write
\[
 Z=\Ric+g,\qquad z=\operatorname{tr}_g Z=s+4,
 \qquad A_0=\operatorname{diag}(-1/6,1/12,1/12),
\]
and let $\mathscr O_0$ be the orthogonal orbit of $A_0$ in
$\operatorname{Sym}^2_0(\Lambda^+)$.  The four-index tensor associated
with a Weyl endomorphism $A$ will be denoted by $\mathcal W(A)$; thus
$|\mathcal W(A)|^2=4|A|^2$.

\subsection{The Cotton source: a weak estimate and an actual metric example}

Use the Cotton convention
\begin{equation}\label{eq:al-cotton}
 C_{ijk}=\nabla_k Z_{ij}-\nabla_j Z_{ik}
 -\frac16\bigl((\nabla_k z)g_{ij}-(\nabla_j z)g_{ik}\bigr).
\end{equation}
With the corresponding divergence sign,
$\operatorname{div}\mathcal W(W^+)=-C^+/2$.  The sign will not matter
in norm estimates.  In particular, the integrated Weyl identity, in these
norms, is
\begin{equation}\label{eq:al-weyl-integral}
 0=\int_M\left\{|\nabla W^+|^2
 -|\operatorname{div}\mathcal W(W^+)|^2
 +\frac{s}{2}|W^+|^2-18\det W^+\right\}.
\end{equation}
This is the usual integrated Weyl--Weitzenb\"ock identity: dividing its
four-index form by four gives the gradient term above, while
$W^+_{ijkl}W^+_{ijpq}W^+_{klpq}=24\det W^+$.  Consequently integration
does not discard the derivative of Ricci; it retains a squared Cotton
term.

For a tensor bundle on a closed metric four-manifold, set
\[
 \|\phi\|_{H^1,\mathrm{av}}^2
 =\frac1V\int_M(|\phi|^2+|\nabla\phi|^2),\qquad
 \|T\|_{H^{-1},\mathrm{av}}
 =\sup_{\|\phi\|_{H^1,\mathrm{av}}\le1}
 \left|\frac1V\langle T,\phi\rangle\right|.
\]
The distribution pairing in this definition is unnormalized before the
factor $V^{-1}$.  These norms belong to the same metric; their definition
does not identify spaces belonging to different manifolds.

\begin{lemma}[Multiplication of the weak Cotton source]
\label{lem:al-weak-source}
Let $a$ be a bounded tensor with $\nabla a\in L^2$ on a closed smooth
four-manifold.  For any fixed tensor contraction of $a$ with $C$,
\begin{equation}\label{eq:al-multiplied-source}
 \|a*C\|_{H^{-1},\mathrm{av}}
 \le c\|Z\|_\infty
 \left(\|a\|_\infty+\|\nabla a\|_{2,\mathrm{av}}\right).
\end{equation}
The constant depends only on the contraction and dimension.  Also,
\begin{equation}\label{eq:al-unmultiplied-source}
 \|\operatorname{div}\mathcal W(W^+)\|_{H^{-1},\mathrm{av}}
 \le c\|Z\|_{2,\mathrm{av}}.
\end{equation}
\end{lemma}

\begin{proof}
All metric contractions and the Hodge projections are covariantly
parallel.  Insert \eqref{eq:al-cotton} and move its derivative onto
$a*\phi$.  The absolute value of the normalized pairing is bounded by
\[
 c\|Z\|_\infty\left(
 \|a\|_\infty\|\nabla\phi\|_{2,\mathrm{av}}
 +\|\nabla a\|_{2,\mathrm{av}}\|\phi\|_{2,\mathrm{av}}\right).
\]
Here normalized Cauchy--Schwarz also bounds the normalized $L^1$ norm
of $\nabla\phi$ by its normalized $L^2$ norm.  Taking the supremum proves
\eqref{eq:al-multiplied-source}.  With no variable multiplier, use
$\|Z\|_2\|\nabla\phi\|_2$ instead.  Approximation, or the weak product
rule directly, proves the assertion for bounded $W^{1,2}$ tensors.
\end{proof}

\begin{proposition}[A bounded-curvature metric sequence with large Cotton energy]
\label{prop:al-cotton-example}
There are smooth metrics $g_j$ on one closed four-manifold such that
\[
 \|\Ric_{g_j}+g_j\|_\infty\longrightarrow0,
 \qquad \sup_j\|\Rm_{g_j}\|_\infty<\infty,
 \qquad \int_M|C^+_{g_j}|^2\longrightarrow\infty.
\]
The sequence may be chosen to converge in $C^1$ to a fixed negative
Einstein metric.
\end{proposition}

\begin{proof}
Take a closed negative Einstein four-metric $g_0$ with $\Ric_{g_0}=-g_0$;
a product of two closed curvature-minus-one surfaces is an example.
On a fixed coordinate patch choose a smooth phase $\theta$ with
$d\theta\ne0$.  Choose a smooth symmetric tensor $q$ satisfying
\[
 \operatorname{tr}_{g_0}q=0,
 \qquad q(\nabla^{g_0}\theta,\cdot)=0,
\]
whose Cotton principal symbol is nonzero on a smaller patch.  Such a
choice is explicit: in an oriented orthonormal frame with
$d\theta=|d\theta|e^1$, use $q=\operatorname{diag}(0,1,-1,0)$ and
extend this frame smoothly after shrinking the patch.  For the index
$i=2$, the skew derivative symbol in \eqref{eq:al-cotton} contains a
nonzero multiple of $e^1\wedge e^2$.  Its self-dual projection is nonzero.
Choose a compactly supported cutoff $\chi$ equal to one on a nonempty
subpatch, and put
\[
 g_{\varepsilon,k}
 =g_0+\frac{\varepsilon}{k^2}\chi\sin(k\theta)q.
\]
These are positive definite for $0<\varepsilon\le\varepsilon_0$ and
large integers $k$, and agree with $g_0$ outside the patch.

The principal symbol of linearized Ricci on a tensor transverse and
trace-free relative to $d\theta$ is $-\tfrac12|d\theta|^2\partial_\theta^2$.
More explicitly, its second-derivative terms are
\[
 \frac12g_0^{ab}
 (\partial_a\partial_i h_{bj}+\partial_a\partial_j h_{bi}
 -\partial_a\partial_b h_{ij}-\partial_i\partial_j h_{ab}),
\]
and the first, second, and fourth principal terms vanish for the chosen
$q$.  It follows, uniformly on the fixed patch, that
\begin{equation}\label{eq:al-ricci-expansion}
 \Ric_{g_{\varepsilon,k}}+g_{\varepsilon,k}
 =\frac\varepsilon2|d\theta|^2\chi\sin(k\theta)q
 +O\left(\frac\varepsilon k+\frac{\varepsilon^2}{k^2}\right).
\end{equation}
Indeed $h=O(\varepsilon k^{-2})$, $\partial h=O(\varepsilon k^{-1})$,
and $\partial^2h=O(\varepsilon)$; the nonlinear second-order terms are
$h*\partial^2h$ and $\partial h*\partial h$.  The same expansion after
one differentiation has remainder $O(\varepsilon+\varepsilon^2/k)$.
The leading scalar term vanishes because $q$ is trace-free.  Thus
\[
 C^+_{g_{\varepsilon,k}}
 =\varepsilon k\cos(k\theta)\,\mathcal T
 +O(\varepsilon+\varepsilon^2/k),
\]
where the fixed smooth tensor $\mathcal T$ is nonzero on a subpatch.
Integration in a coordinate having $\theta$ as one coordinate gives
\[
 \int\cos^2(k\theta)|\mathcal T|^2\,d\mu_{g_0}
 \longrightarrow\frac12\int|\mathcal T|^2\,d\mu_{g_0}>0.
\]
The volume forms of the perturbed metrics converge uniformly to that of
$g_0$.  Consequently
$\int|C^+_{g_{\varepsilon,k}}|^2\ge c\varepsilon^2k^2$ for sufficiently
large $k$, uniformly for $0<\varepsilon\le\varepsilon_0$.
The coordinate curvature formula and the derivative bounds on $h$ give
$\|\Rm_{g_{\varepsilon,k}}\|_\infty\le C$.  Finally take
$\varepsilon_j=j^{-1}$ and $k_j=j^2$.  Equation
\eqref{eq:al-ricci-expansion} proves the residual assertion, while the
Cotton energy is at least $cj^2$.
\end{proof}

The example does not approach the forbidden Weyl orbit: its limiting
metric is an ordinary negative Einstein product.  It disproves a bound
for Cotton energy from small Ricci residual and bounded curvature, not
the calibrated obstruction theorem.

\subsection{Two critical concentration tests}

For $A\in\operatorname{Sym}^2_0(\mathbb R^3)$ and $s\in\mathbb R$, put
\[
 f=\frac{2s}{3}-2\sqrt{\frac23}|A|,
 \quad b=\frac{2s}{3}+2\lambda_{\min}(A),
 \quad G=f^2-b_-^2,\qquad b_-:=\min(b,0).
\]
Since $f\le b$, $G\ge0$.  The scalar--spectral density is
\begin{equation}\label{eq:al-density}
 \mathscr D=|Z|^2+(s+4)^2+(s+8\sqrt6|A|)^2+G.
\end{equation}
For $Z=0$, $s=-4$, and $A=A_0$, this density is zero.

\begin{proposition}[Quadratic control does not control the cubic invariant]
\label{prop:al-cubic-concentration}
On a fixed unit-volume Euclidean domain there are smooth algebraic Weyl
fields $A_j$ with
\[
 \int\mathscr D_j\longrightarrow0,
 \qquad \|A_j-A_0\|_2\longrightarrow0,
 \qquad \int(\det A_j-\det A_0)\longrightarrow+\infty.
\]
Here the Ricci and scalar entries of the algebraic curvature tensor are
fixed to be $-g$ and $-4$.  No metric realizing these curvature fields is
asserted.
\end{proposition}

\begin{proof}
Let $\eta\ge0$ be a nonzero smooth bump supported in the unit ball, and
take $t_j\to\infty$, $r_j=t_j^{-5/8}$, and
\[
 B=\operatorname{diag}(2,-1,-1),\qquad
 A_j=A_0+t_j\eta(x/r_j)B.
\]
For large $j$ the support lies inside the domain.  Its volume is
$O(r_j^4)=O(t_j^{-5/2})$.  Outside that support the density vanishes;
inside, every summand in \eqref{eq:al-density} is bounded by
$C(1+t_j^2\eta(x/r_j)^2)$.  Thus $\int\mathscr D_j=O(t_j^{-1/2})$.
Writing $c_m=\int\eta^m$, one has exactly
\[
 \|A_j-A_0\|_2^2=6c_2t_j^2r_j^4=6c_2t_j^{-1/2}.
\]
Since $A_0=-B/12$ and $\det B=2$, expansion of the determinant gives
\[
 \int(\det A_j-\det A_0)
 =r_j^4\left(2c_3t_j^3-\frac{c_2}2t_j^2
                         +\frac{c_1}{24}t_j\right)
 =2c_3t_j^{1/2}+o(1).
\]
The usual curvature-operator block decomposition realizes these entries
as algebraic curvature tensors pointwise.  It imposes no differential
Bianchi identity, which is why this construction makes no geometric
realizability claim.
\end{proof}

\begin{proposition}[A small $L^1$ commutator with a non-small $H^{-1}$ norm]
\label{prop:al-commutator}
On a fixed Euclidean four-ball there are smooth functions $a_r,R_r$,
compactly supported in a ball of radius $O(r)$, such that
\[
 \|a_r\|_\infty\le C,\qquad
 \|\nabla a_r\|_2+\|R_r\|_2=O(r^{1/2}),\qquad
 \|R_r\partial_1a_r\|_1=O(r),
\]
but $\liminf_{r\downarrow0}\|R_r\partial_1a_r\|_{H^{-1}}>0$.
Here $H^{-1}$ is the dual of the inhomogeneous $H^1_0$ norm on the fixed
ball.
\end{proposition}

\begin{proof}
Choose a nonzero smooth compactly supported bump $\chi$ and a smooth
compactly supported $\zeta$ which is positive on its support.  Set
\[
 a_r(x)=\chi(x/r)\sin(x_1/r^{3/2}),\quad
 R_r(x)=r^{-3/2}\chi(x/r)\cos(x_1/r^{3/2}),\quad
 \phi_r(x)=r^{-1}\zeta(x/r).
\]
The change of variables $x=ry$ gives
$\|\phi_r\|_2=O(r)$, $\|\nabla\phi_r\|_2=O(1)$, and the claimed
bounds on $a_r,R_r$.  The main derivative of $a_r$ has size $r^{-3/2}$
on volume $O(r^4)$; the cutoff derivative has size $r^{-1}$.
Cauchy--Schwarz therefore gives the $L^1$ bound.  On the other hand,
\[
 \int R_r(\partial_1a_r)\phi_r
 =\int\chi(y)^2\zeta(y)\cos^2(y_1/\sqrt r)\,dy+O(r^{1/2})
 \longrightarrow\frac12\int\chi^2\zeta>0.
\]
The uniformly bounded $H^1_0$ test norms prove the lower bound.
\end{proof}

The angle $a_r$ can be used in a unit-vector rotation
$(\cos a_r,\sin a_r)$, or in a unit self-dual-form rotation in a fixed
parallel frame.  Its amplitude is then exactly constant, and its
covariant derivative and exterior derivative tend to zero in $L^2$.
The example still imposes neither the metric curvature equation nor a
monopole equation.  It excludes the inference
$R\nabla a\to0$ in $H^{-1}$ from the stated norms alone.  In contrast,
\eqref{eq:al-multiplied-source} exploits the fact that Cotton is a
derivative of the \emph{bounded} Ricci residual.

\subsection{Global rank-one alignment without a global eigenvalue gap}

The following statement uses the actual monopole supply in
Section~\ref{sec:early-walsh}: every specified sign class admits the
rescaled monopole solutions required there.  A collection of formal
cohomology classes does not suffice.  All norms in this subsection are
unnormalized.

\begin{lemma}[Good-set alignment and bad-set mass]
\label{lem:al-global-phase}
Assume $\delta=\|g^{-1}\Ric+I\|_{\infty,\mathrm{op}}\le1/64$.
Let $x_1,\ldots,x_N$ be harmonic self-dual two-forms such that every
$H_\epsilon=\sum_r\epsilon_rx_r$, $\epsilon\in\{-1,1\}^N$, is the
harmonic self-dual representative of a specified actual monopole class
with the solvability hypothesis just stated.  Put
\[
 \begin{gathered}
 q=-s+\sqrt6|W^+|,\quad
 Q=\int q^2,\quad \kappa=\frac{Q}{72\pi^2},\quad
 \Delta_\epsilon=Q-72\pi^2\|H_\epsilon\|_2^2,\\
 L=\frac{q}{6\sqrt2\pi},\quad c_* =\frac3{4\sqrt2\pi},\quad
 d_{\rm cal}=\int\mathscr D,\quad
 \overline\Delta=2^{-N}\sum_\epsilon\Delta_\epsilon,\quad
 B_0=\frac{\overline\Delta+37d_{\rm cal}}{\pi^2},
 \end{gathered}
\]
where $\mathscr D$ is \eqref{eq:al-density} with $A=W^+$.
There is a measurable orthogonal projection $P$, of rank one on
$U=\{\mathscr D\le1/64\}$ and zero on its complement, such that
\begin{equation}\label{eq:al-global-projection}
 \|(I-P)H_\epsilon\|_2^2
 \le\frac{\Delta_\epsilon+37d_{\rm cal}}{\pi^2},
 \qquad \sum_r\|(I-P)x_r\|_2^2\le B_0.
\end{equation}
Locally write $Px_r=\alpha_rv$ for a unit section $v$ of its image, and
set $\alpha_r=0$ off $U$.  The sign-independent products obey
\begin{align}
 \left\|\sum_r\alpha_r^2-L^2\right\|_1
 &\le2\sqrt{\frac{\kappa\overline\Delta}{18\pi^2}}+B_0,
 \label{eq:al-phase-mass}\\
 \|\alpha_p\alpha_q\|_1
 &\le\sqrt{\frac{\kappa\overline\Delta}{18\pi^2}}+\frac{B_0}{2}
 \quad(p\ne q).
 \label{eq:al-phase-products}
\end{align}
One may replace $L^2$ by $c_*^2$ in \eqref{eq:al-phase-mass} at the
additional cost
\begin{equation}\label{eq:al-constant-phase-target}
 (\sqrt\kappa+c_*\sqrt V)
       \sqrt{\frac{41d_{\rm cal}}{2304\pi^2}}.
\end{equation}
The constants are independent of the size of the sign cube.
\end{lemma}

\begin{proof}
The amplitude and spectral-deficit estimates of
Section~\ref{sec:early-walsh}, in precisely this normalization, give
$\Delta_\epsilon\ge0$ and
\begin{equation}\label{eq:al-phase-input}
 \||H_\epsilon|-L\|_2^2\le\frac{\Delta_\epsilon}{18\pi^2}.
\end{equation}
On a measurable region where the lowest eigenvalue gap is at least
$\gamma$ and $q\le F^2$, they also give
\[
 \int|\Pi^\perp H_\epsilon|^2
 \le\left(\frac1{36\pi^2}
                  +\frac{F^2}{216\gamma\pi^2}\right)\Delta_\epsilon.
\]
These inequalities follow by smooth positive approximation of
$\sqrt q$; the harmonic forms are fixed during approximation, so the
displayed inequalities have no remaining allowance.

Write $a=\sqrt{2/3}|W^+|$, $w=\lambda_{\min}(W^+)$, and
$\zeta=w+a\ge0$.  On $U$, $|t|\le1/8$, where $t=s+24a$, and
$|z|=|s+4|\le1/16$.  Hence
\[
 |a-1/6|\le1/128,\qquad \frac{61}{384}\le a\le\frac{67}{384}.
\]
Since $s<-3$ and $w\le0$, the quantity $b$ is negative and
$G=4\zeta(-2s/3+2a-\zeta)\ge8\zeta$.
In particular $\zeta\le1/512$ on $U$.  Trace-freeness gives
$\lambda_{2,3}=-w/2\mp\tfrac12\sqrt{2|W^+|^2-3w^2}$, so
\[
 \lambda_2-w
 \ge\frac32\left(\frac{61}{384}-\frac1{512}\right)
       -\frac12\sqrt{\frac{6\cdot67}{384\cdot512}}>\frac18.
\]
Also $q=-s+3a\le587/128<5$.  The spectral estimate thus applies on
$U$ with $\gamma=1/8$ and $F^2=5$, giving
\[
 \int_U|\Pi^\perp H_\epsilon|^2
       \le\frac{23\Delta_\epsilon}{108\pi^2}.
\]
Take $P=\Pi$ there and $P=0$ on $B=M\setminus U$.

Chebyshev's inequality gives $\Vol(B)\le64d_{\rm cal}$.  More
importantly, the exact scalar identity $q-9/2=(t-9z)/8$ gives
\[
 \|q-9/2\|_2^2\le\frac{41}{32}d_{\rm cal},\qquad
 \||H_\epsilon|-c_*\|_2^2
 \le\frac{\Delta_\epsilon}{9\pi^2}
                +\frac{41d_{\rm cal}}{1152\pi^2}.
\]
Therefore the entire harmonic mass on the bad set satisfies
\[
 \int_B|H_\epsilon|^2
 \le2\||H_\epsilon|-c_*\|_2^2+2c_*^2\Vol(B)
 \le\frac1{\pi^2}\left\{\frac29\Delta_\epsilon
                       +\left(36+\frac{41}{576}\right)d_{\rm cal}\right\}.
\]
Adding the good-set estimate proves the first part of
\eqref{eq:al-global-projection}, since $23/108+2/9=47/108<1$.
This step needs neither a maximum bound for $H_\epsilon$ nor a
Sobolev inequality.

Write $y_r=(I-P)x_r$ and average over independent signs.  Walsh
orthogonality yields
$\sum_r\|y_r\|_2^2=2^{-N}\sum_\epsilon\|(I-P)H_\epsilon\|_2^2$,
proving the second assertion.  Since $\|H_\epsilon\|_2\le\sqrt\kappa$
and $\|L\|_2=\sqrt\kappa$, \eqref{eq:al-phase-input} implies
\[
 \||H_\epsilon|^2-L^2\|_1
 \le2\sqrt{\frac{\kappa\Delta_\epsilon}{18\pi^2}}.
\]
The constant and pairwise sign coefficients of $|H_\epsilon|^2$ are
$\sum_r|x_r|^2$ and $2\langle x_p,x_q\rangle$, respectively.
Taking their $L^1$ norms and using the sign average of
$\sqrt{\Delta_\epsilon}$ at most $\sqrt{\overline\Delta}$ gives
the corresponding bounds $2\sqrt{\kappa\overline\Delta/(18\pi^2)}$
and $\sqrt{\kappa\overline\Delta/(18\pi^2)}$.
The decompositions $x_r=Px_r+y_r$ are orthogonal pointwise.
Subtract $\sum_r|y_r|^2$ for \eqref{eq:al-phase-mass}, and
$\langle y_p,y_q\rangle$ for \eqref{eq:al-phase-products}, using
$\|y_p\|_2\|y_q\|_2\le B_0/2$.
Finally $\|L-c_*\|_2^2\le41d_{\rm cal}/(2304\pi^2)$, and
Cauchy--Schwarz applied to $(L-c_*)(L+c_*)$ proves
\eqref{eq:al-constant-phase-target}.
\end{proof}

For example, the volume on which $|\alpha_p|,|\alpha_q|\ge h>0$
is at most $h^{-2}$ times the right side of
\eqref{eq:al-phase-products}.  These are global mass and overlap
estimates.  The projection is only measurable: its derivative is never
taken, and no global unit eigenvector is chosen.  Thus the lemma does
not assert a common controlled coordinate patch or exclude separated
phases connected through regions of small measure.

\subsection{The bounded Weyl model has an exact Hodge cancellation}

For a smooth self-dual two-form $\psi$, write $\rho=|\psi|$ and define
\begin{equation}\label{eq:al-polynomial-model}
 A_\psi=\frac{\rho^2}{24}I-\frac18\psi\otimes\psi,
 \qquad B_\psi=\operatorname{div}\mathcal W(A_\psi),
 \qquad K_\psi=\frac{s}{3}\rho^2-2\langle W^+\psi,\psi\rangle.
\end{equation}
The operator is smooth even at zeros of $\psi$.

\begin{lemma}[Exact model identity]\label{lem:al-hodge-cancellation}
For every smooth self-dual two-form on a closed oriented four-manifold,
\begin{equation}\label{eq:al-model-identity}
 \int_M\left\{|\nabla A_\psi|^2-|B_\psi|^2
       +\frac{s}{2}|A_\psi|^2
       -6\operatorname{tr}(W^+A_\psi^2)\right\}=0.
\end{equation}
This identity is a weighted two-form Hodge identity and imposes no
additional equation on $\psi$.
\end{lemma}

\begin{proof}
The pointwise first-derivative calculation in the stated tensor norms is
\begin{align}
 |\nabla A_\psi|^2-|B_\psi|^2
 &=\frac{\rho^2}{64}(|\nabla\psi|^2-2|d\psi|^2)
   +\frac{\rho^2}{32}|d\rho|^2
   -\frac{\rho}{16}\langle\iota_{\nabla\rho}\psi,*d\psi\rangle.
 \label{eq:al-first-jet}
\end{align}
Here is an explicit way to perform the contraction.  At a point where
$\rho>0$, choose an orthonormal basis $e_a=\omega_a/\sqrt2$ of
$\Lambda^+$, with
\[
 \omega_1=dx^{01}+dx^{23},\quad
 \omega_2=dx^{02}+dx^{31},\quad
 \omega_3=dx^{03}+dx^{12},
\]
and put $\psi=\rho e_1$ and
$\nabla_i\psi=a_i e_1+b_i e_2+c_i e_3$.  Thus $a=d\rho$.
In this basis,
\[
 \nabla_iA_\psi=\rho
 \begin{pmatrix}
 -a_i/6&-b_i/8&-c_i/8\\
 -b_i/8&a_i/12&0\\
 -c_i/8&0&a_i/12
 \end{pmatrix},\qquad
 (B_\psi)_{jkl}
 =\frac12\sum_{i,a,b}(\nabla_iA_\psi)_{ab}
                         (\omega_a)_{ij}(\omega_b)_{kl}.
\]
Use $\sum_{kl}(\omega_a)_{kl}(\omega_b)_{kl}=4\delta_{ab}$,
$d\psi=\sum_i dx^i\wedge\nabla_i\psi$, and
$*(\xi\wedge\psi)=\iota_{\xi^\sharp}\psi$.
Squaring the two displayed expressions and collecting the radial,
transverse, and mixed terms gives \eqref{eq:al-first-jet}.  In particular
the same matrix gives
\begin{equation}\label{eq:al-model-gradient}
 |\nabla A_\psi|^2
 =\frac{\rho^2}{32}|\nabla\psi|^2
  +\frac{\rho^2}{96}|d\rho|^2.
\end{equation}
All weighted expressions extend across zeros: $\rho\,d\rho
=\tfrac12d(\rho^2)$ and $\rho^2|d\rho|^2=|d(\rho^2)|^2/4$.

The eigenvalues of $A_\psi$ are
$-\rho^2/12,\rho^2/24,\rho^2/24$.  Therefore
\[
 |A_\psi|^2=\frac{\rho^4}{96},\qquad
 \operatorname{tr}(W^+A_\psi^2)
 =\frac{\rho^2}{192}\langle W^+\psi,\psi\rangle,
\]
so the last two terms of \eqref{eq:al-model-identity} equal
$\rho^2K_\psi/64$.

Finally test
$\Delta_H\psi=\nabla^*\nabla\psi+(s/3)\psi-2W^+\psi$
against $w\psi$, where $w=\rho^2$.  Self-duality gives
\[
 \langle d\psi,d(w\psi)\rangle
 +\langle d^*\psi,d^*(w\psi)\rangle
 =2w|d\psi|^2+2\langle*d\psi,\iota_{\nabla w}\psi\rangle.
\]
The rough-Laplacian pairing is
$w|\nabla\psi|^2+|dw|^2/2$.  Hence
\begin{equation}\label{eq:al-weighted-hodge}
 \int\rho^2(|\nabla\psi|^2+K_\psi-2|d\psi|^2)
 =4\int\rho\langle\iota_{\nabla\rho}\psi,*d\psi\rangle
  -2\int\rho^2|d\rho|^2.
\end{equation}
Substituting this equality into \eqref{eq:al-first-jet} cancels both
remaining terms and proves \eqref{eq:al-model-identity}.
\end{proof}

\begin{corollary}[A weakly vanishing divergence can retain energy]
\label{cor:al-surviving-energy}
Suppose $(M_j,g_j,\psi_j)$ is a sequence of closed oriented metric
four-manifolds with smooth self-dual forms, and, in normalized norms,
\[
 \sup_j\|\rho_j\|_\infty<\infty,\qquad
 \rho_j\longrightarrow\sqrt2\text{ in }L^2,\qquad
 \|d\rho_j\|_2+\|d\psi_j\|_2+\|Z_j\|_2
       +\|W^+_j-A_{\psi_j}\|_2\longrightarrow0.
\]
Then the normalized energies satisfy
\[
 \|\nabla\psi_j\|_2^2\longrightarrow2,\quad
 \frac1{V_j}\int\rho_j^2|\nabla\psi_j|^2\longrightarrow4,
 \quad \|\nabla A_{\psi_j}\|_2^2\longrightarrow\frac18,
 \quad \|B_{\psi_j}\|_2^2\longrightarrow\frac1{16},
\]
although $\|B_{\psi_j}\|_{H^{-1},\mathrm{av}}\to0$.
\end{corollary}

\begin{proof}
Put $R_j=W^+_j-A_{\psi_j}$.  Since
$A_\psi\psi=-\rho^2\psi/12$,
\[
 K_{\psi_j}=(s_j/3)\rho_j^2+\rho_j^4/6
                   -2\langle R_j\psi_j,\psi_j\rangle.
\]
Bounded amplitudes and their normalized $L^2$ convergence imply convergence
of every finite amplitude moment.  Also $s_j\to-4$ in normalized $L^2$.
Consequently the mean of $K_{\psi_j}$ tends to $-2$, and the mean of
$\rho_j^2K_{\psi_j}$ tends to $-4$.  The unweighted Hodge identity gives
the first limit.  In \eqref{eq:al-weighted-hodge} the two right-hand terms
tend to zero by Cauchy--Schwarz, giving the second.  Equation
\eqref{eq:al-model-gradient} then gives $1/8$, and
\eqref{eq:al-first-jet} gives $1/16$.  No uniform integrability of the
gradient energy has been used to interchange a weight and a limit.
Finally
\[
 B_{\psi_j}=\operatorname{div}\mathcal W(W^+_j)
             -\operatorname{div}\mathcal W(R_j),
\]
and integration by parts, together with
\eqref{eq:al-unmultiplied-source}, bounds its negative norm by
$c(\|Z_j\|_{2,\mathrm{av}}+\|R_j\|_{2,\mathrm{av}})$.
\end{proof}

This is a conditional implication, not the construction of such a
sequence.  It identifies the information lost by replacing a weak Bianchi
source with a small $L^2$ quantity.  The covariant Hessian self-test of
Section~\ref{sec:intrinsic-calibration} supplies information beyond these
first-order identities.

\subsection{The single-chirality symbol and controlled-chart transfer}

\begin{lemma}[The Weyl-half divergence symbol]\label{lem:al-symbol}
In oriented Euclidean four-space let
\[
 \mathcal W(A)_{ijkl}
 =\frac12\sum_{a,b}A_{ab}(\omega_a)_{ij}(\omega_b)_{kl},
 \qquad DA=\partial^i\mathcal W(A)_{ijkl},
\]
using the forms $\omega_a$ in the proof of
Lemma~\ref{lem:al-hodge-cancellation}.  For every covector $\xi$,
\begin{equation}\label{eq:al-symbol}
 |\sigma_D(\xi)A|^2=|\xi|^2|A|^2.
\end{equation}
On the constant-spectrum orbit $\mathscr O_0$, the complete covariant
first jet satisfies
\begin{equation}\label{eq:al-orbit-jet}
 \frac12|\nabla A|^2
 \le|\operatorname{div}\mathcal W(A)|^2
 \le\frac32|\nabla A|^2.
\end{equation}
\end{lemma}

\begin{proof}
By orthogonal covariance it suffices to take $\xi=dx^0$.  The vectors
$((\omega_a)_{0j})_j$ are orthonormal, whereas the squared tensor inner
products of the $\omega_b$ are $4\delta_{bd}$.  Contraction in the formula
for $\mathcal W$ gives \eqref{eq:al-symbol}.

At a diagonal point of the orbit, each derivative has the form
\[
 H_i=\begin{pmatrix}0&u_i&v_i\\u_i&0&0\\v_i&0&0\end{pmatrix}.
\]
The divergence tensor is
\[
 \frac12\left\{
 (\iota_u\omega_1)\otimes\omega_2
 +(\iota_v\omega_1)\otimes\omega_3
 +(\iota_u\omega_2+\iota_v\omega_3)\otimes\omega_1\right\}.
\]
Its squared norm is
$2(|u|^2+|v|^2)+2\langle\iota_u\omega_2,\iota_v\omega_3\rangle$.
Each map $u\mapsto\iota_u\omega_a$ is an isometry.  Bounding the cross
term by $|u||v|$ places this expression between
$|u|^2+|v|^2$ and $3(|u|^2+|v|^2)$.  Since
$\sum_i|H_i|^2=2(|u|^2+|v|^2)$, this proves
\eqref{eq:al-orbit-jet}.
\end{proof}

For example, on a flat torus Fourier transformation gives the exact
identity $\|A-\overline A\|_2^2=\|DA\|_{\dot H^{-1}}^2$, omitting the
zero Fourier mode.  Nevertheless, if $W^+=A+R$, a small $L^2$ norm of
$R$ controls its divergence only weakly.  Expressing that divergence in a
varying eigenframe introduces precisely a product of the type in
Proposition~\ref{prop:al-commutator}.

\begin{proposition}[Strong curvature transfer on controlled charts]
\label{prop:al-chart-transfer}
Let $U$ be a fixed coordinate ball and let $g_j$ be smooth metrics,
uniformly elliptic and bounded in $C^{1,\alpha}(U)$ for some $\alpha>0$.
Suppose $g_j\to g_\infty$ in $C^1(U)$ and
\[
 \|\Ric(g_j)+g_j\|_{L^2(U)}\longrightarrow0.
\]
After the algebraic metric-dependent identifications of the two self-dual
bundles, $W^+(g_j)$ converges strongly in $L^2(U')$ for every $U'\Subset U$.
The limit metric is a smooth negative Einstein metric in its harmonic
coordinates.  A separate $L^2$ curvature bound is unnecessary.
\end{proposition}

\begin{proof}
All norms in this proof are fixed coordinate norms; uniform ellipticity
makes them equivalent to the metric norms.  In algebraically selected
orthonormal frames the divergence operator is
$D_j=P_j^a\partial_a+Q_j$, with $P_j\to P_\infty$ in $C^1$ and
$Q_j\to Q_\infty$ uniformly.  Integration by parts gives, on fixed
interior supports,
\begin{equation}\label{eq:al-operator-convergence}
 \|(D_j-D_\infty)u\|_{H^{-1}}=o(1)\|u\|_2,
 \qquad \|D_jW^+(g_j)\|_{H^{-1}}
       \le C\|\Ric(g_j)+g_j\|_2.
\end{equation}
The second estimate is \eqref{eq:al-cotton} integrated against an $H^1_0$
test tensor.  The derivatives of the volume density and tensor
coefficients are bounded by the assumed coordinate control.

The injective symbol in \eqref{eq:al-symbol} gives the negative-order
interior estimate
\begin{equation}\label{eq:al-negative-elliptic}
 \|u\|_2\le C\bigl(\|D_\infty u\|_{H^{-1}}
                           +\|u\|_{H^{-1}}\bigr)
\end{equation}
for fields supported in a fixed compact subset of $U$.  For constant
coefficients this follows by Fourier transformation; the second term
accounts for low frequencies.  To obtain the variable-coefficient
estimate, use finitely many small coordinate balls and freeze the
principal coefficients.  Writing
$(P(x)-P(x_0))\partial u
=\partial((P(x)-P(x_0))u)-\partial P(x)u$,
the first term costs at most the oscillation of $P$ times $\|u\|_2$.
The remaining zeroth-order maps from $L^2$ to $H^{-1}$ are compact,
including the localization commutators.  Each satisfies
\[
 \|Ku\|_{H^{-1}}\le\eta\|u\|_2+C_\eta\|u\|_{H^{-1}}.
\]
Indeed a contrary normalized sequence would converge weakly to zero
in $L^2$ and strongly to zero in $H^{-1}$, while its compact image did
not converge to zero.  Absorb both the coefficient oscillation and these
$\eta$ terms to prove \eqref{eq:al-negative-elliptic}.

Set $A_j=W^+(g_j)$.  The curvature formula
$\partial\Gamma_j+\Gamma_j*\Gamma_j$ and $C^1$ convergence of the
metrics show first that $A_j$ converges strongly in $H^{-1}$ locally,
to a distribution $A_\infty$.  Contraction coefficients and the frame
identifications converge in $C^1$, so multiplication by them is continuous
on $H^{-1}$.  This argument does not presuppose an $L^2$ curvature bound.
For an interior cutoff $\chi$,
\[
 D_\infty(\chi A_j)
 =(D_\infty-D_j)(\chi A_j)+\chi D_jA_j+[D_j,\chi]A_j.
\]
Apply \eqref{eq:al-negative-elliptic}.  The first term is
$o(1)\|\chi A_j\|_2$ and is absorbed.  The second tends to zero by
\eqref{eq:al-operator-convergence}.  The commutator is multiplication by
$P_j^a\partial_a\chi$, bounded in $C^1$, so its negative norm is bounded
by the already established local $H^{-1}$ control of curvature.
Thus $\|\chi A_j\|_2$ is uniformly bounded.  Repeat the estimate on
$\chi(A_j-A_k)$.  Operator differences now tend to zero using this
$L^2$ bound; the sources tend to zero; and the commutators converge in
$H^{-1}$ using the $C^1$ convergence of their coefficients.  The sequence
is therefore Cauchy in $L^2$ on the region where $\chi=1$.

Distributional convergence identifies the limit with $W^+(g_\infty)$
and gives $\Ric(g_\infty)=-g_\infty$.  The $C^{1,\alpha}$ bound passes
to the limit.  Construct harmonic coordinates for this limit metric on
a smaller ball by solving its scalar Dirichlet problems.  Shrinking the
ball gives an invertible coordinate differential.  In those coordinates
the weak Ricci equation is a uniformly elliptic system, with principal
part $-\tfrac12g_\infty^{ab}\partial_a\partial_b g_\infty$ and lower-order
terms quadratic in first derivatives.  The divergence formulation first
gives interior $W^{2,p}$ regularity, and Schauder estimates then bootstrap
to smoothness.  Only the limit is put in harmonic coordinates; a common
harmonic gauge for all $g_j$ was not assumed.
\end{proof}

\subsection{A geometric gap under bounded full curvature}

\begin{theorem}[A collapse-tolerant bounded-curvature gap]
\label{thm:al-bounded-curvature-gap}
For every finite $K\ge0$ there is $c(K)>0$ such that every closed oriented
smooth four-manifold with $|\Rm|\le K$ satisfies
\begin{equation}\label{eq:al-bounded-curvature-gap}
 \frac1V\int_M\left(|\Ric+g|^2+
                     \operatorname{dist}(W^+,\mathscr O_0)^2\right)
 \ge c(K).
\end{equation}
No volume, diameter, injectivity-radius, Cotton, or Ricci-derivative bound
is imposed.  In particular, if
$\delta=\|g^{-1}\Ric+I\|_{\infty,\mathrm{op}}$, then
$\delta^2+V^{-1}\int\operatorname{dist}(W^+,\mathscr O_0)^2\ge c(K)/4$.
\end{theorem}

\begin{proof}
Choose $\rho=1/(100\sqrt{1+K})$.  The exponential map at any point $p$
has no conjugate points on the tangent ball of radius $\rho$.  On this
ball its pullback metric $h$ satisfies
$Q^{-1}I\le h\le QI$, with a universal $Q$, by Rauch comparison.
The map need not be injective into $M$.  In the source, however,
Gauss's lemma gives $d_h(0,v)=|v|$.  Every sufficiently small interior
source ball satisfies
\[
 \Vol_h B_t(x)\ge Q^{-4}\omega_4t^4,
\]
because it contains the Euclidean ball of radius $t/\sqrt Q$ and the
volume density is at least $Q^{-2}$.  Together with the curvature bound
on a buffered larger ball, the interior bounded-geometry harmonic-radius
estimate used in Section~\ref{sec:early-local-rigidity} supplies uniform
$C^{1,\alpha}$ harmonic charts on a smaller fixed ball.  This application
uses source noncollapse and the buffer, not injectivity radius on $M$.

Fix $s>0$ sufficiently smaller than $\rho$, so that the closed source
$3s$-ball lies in this development.  The buffered local-isometry transfer
lemma of Section~\ref{sec:early-local-rigidity} states, for $H\ge0$,
\begin{equation}\label{eq:al-buffer-transfer}
 \frac1{\Vol_hB_s(0)}\int_{B_s(0)}\exp_p^*H
 \le\frac{\Vol_h B_{3s}(0)}{\Vol_h B_s(0)}
       \frac1{\Vol_g B_s(p)}\int_{B_s(p)}H.
\end{equation}
Its multiplicity argument compares preimages by lifting paths of length
less than $2s$ inside the compact buffer; it assumes neither a covering
map nor a uniform multiplicity bound.  The source volume ratio is bounded
by $3^4Q^8$.

If the theorem were false, choose a sequence with the left side of
\eqref{eq:al-bounded-curvature-gap} tending to zero, and denote its
nonnegative integrand by $H_j$.  Ricci has a uniform lower bound.
Bishop--Gromov comparison and Fubini yield
\[
 \int_M\left(\frac1{\Vol B_s(p)}\int_{B_s(p)}H_j\right)d\mu(p)
 \le D(K,s)\int_M H_j.
\]
Indeed, if $q\in B_s(p)$, then $B_s(q)\subset B_{2s}(p)$, whose volume
is at most $D(K,s)\Vol B_s(p)$; exchange the two integrations.
Choose a center with ball average at most $D(K,s)$ times the global
average.  Equation \eqref{eq:al-buffer-transfer} transfers its vanishing
average to the fixed, noncollapsed source ball.  On a smaller fixed
harmonic chart the metrics have a $C^1$ convergent subsequence and
$\|\Ric+g\|_2\to0$.  Proposition~\ref{prop:al-chart-transfer} gives
strong $L^2$ Weyl convergence and a smooth negative Einstein limit.
Distance to a fixed closed orbit is Lipschitz, so the limit takes values
in $\mathscr O_0$ on a nonempty open set.

An Einstein Weyl tensor has zero divergence.  Inequality
\eqref{eq:al-orbit-jet} makes this constant-spectrum Weyl tensor parallel.
Writing it as $I/12-(n\otimes n)/4$ shows that its simple eigenline is
parallel.  On a smaller simply connected patch its unit self-dual section
therefore defines a parallel complex structure.  The metric is locally
K\"ahler and must satisfy $|W^+|^2=s^2/24=2/3$, whereas the prescribed
orbit has squared norm $1/24$.  This contradiction proves the theorem.
The final assertion follows from $|\Ric+g|^2\le4\delta^2$.
\end{proof}

For clarity, the scalar--spectral density directly controls the distance
used in this theorem when $s\le-3$.  Put
$u=|W^+|$, $\alpha=\sqrt{2/3}\,u$, $\zeta=\lambda_{\min}(W^+)+\alpha$,
and $t=s+8\sqrt6u$.  Then $0\le\zeta\le\alpha$ and
\[
 \operatorname{dist}(W^+,\mathscr O_0)^2
 =\left(u-\frac1{2\sqrt6}\right)^2+\frac\zeta2
 \le\frac{t^2+(s+4)^2}{192}+\frac G{16}.
\]
The equality follows by projecting onto the lowest eigenline.  For the
last inequality use $t-(s+4)=8\sqrt6u-4$ and
$G=4\zeta(-2s/3+2\alpha-\zeta)\ge8\zeta$.
Consequently $|\Ric+g|^2+\operatorname{dist}(W^+,\mathscr O_0)^2
\le\mathscr D$.  The actual monopole density estimate proved in
Section~\ref{sec:early-local-rigidity},
$V^{-1}\int\mathscr D\le436\delta^2+675d$, therefore gives
\[
 \delta^2+d\ge\frac{c(K)}{675}
\]
under its specified monopole and topology hypotheses, with
$d=(T/A-2/3)_+$ and $\delta\le1/4$.

There is also a localized consequence.  If $\delta_j\to0$ and
$V_j^{-1}\int\operatorname{dist}(W_j^+,\mathscr O_0)^2\to0$, then for
every fixed $K,R>0$,
\begin{equation}\label{eq:al-localized-degeneration}
 \frac{\Vol\{p:\sup_{B_R(p)}|\Rm_{g_j}|\le K\}}{V_j}
 \longrightarrow0.
\end{equation}
Otherwise a positive volume proportion of such centers remains.  Choose
the development radius above also less than $R/4$.  Fubini restricted to
that set selects a center with vanishing calibration average, at the cost
of the reciprocal of its positive volume fraction.  The same buffered
development and limiting argument give a contradiction.  Statement
\eqref{eq:al-localized-degeneration} concerns fixed-radius neighborhoods
of concentration, not pointwise blowup at every center.  None of these
controlled-geometry arguments replaces the intrinsic Hessian estimate
when such charts and curvature bounds are unavailable.
